% 原创教学示意；超边保留参与组，右侧为节点—关系二部表示。
\begin{tikzpicture}[x={\dimexpr\linewidth/112\relax},y=1mm,
  font=\figurefont\footnotesize,text=NNInk]
  \draw[draw=NNHidden,fill=NNHidden!7,line width=.8pt,rounded corners=4mm]
    (0,-7)--(0,21)--(19,25)--(37,6)--(37,-7)--cycle;
  \draw[draw=NNOutput,fill=NNOutput!7,line width=.8pt,rounded corners=4mm]
    (22,-24) rectangle (47,9);
  \foreach \i/\xx/\yy in {1/8/16,2/8/0,3/29/0,4/39/-16}{
    \node[nn input,minimum size=6mm] (v\i) at (\xx,\yy) {$\i$};
  }
  \node[text=NNHidden] at (23,17) {$S_a$};
  \node[text=NNOutput] at (28,-17) {$S_b$};
  \draw[nn link] (53,-27)--(53,30);
  \foreach \i/\yy in {1/24,2/8,3/-8,4/-24}{
    \node[nn input,minimum size=6mm] (u\i) at (68,\yy) {$\i$};
  }
  \node[nn operator,minimum width=12mm] (sa) at (104,13) {$S_a$};
  \node[nn operator,draw=NNOutput,fill=NNOutput!8,minimum width=12mm] (sb) at (104,-17) {$S_b$};
  \foreach \i in {1,2,3}{\draw[nn link,draw=NNHidden] (u\i)--(sa);}
  \foreach \i in {3,4}{\draw[nn link,draw=NNOutput] (u\i)--(sb);}
\end{tikzpicture}
